\section{Discussion}

\paragraph{A compliance spectrum, not a binary.}
Our results replace a pass/fail view of policy compliance with a spectrum:
strict refusal (Haiku, Opus, GPT-4o), probing-then-disavowal (Sonnet), and
genuine violation (GPT-4o-mini, Llama).
The probe behavior is itself a form of misalignment between the behavioral and
natural-language channels---a tool-action analogue of unfaithful
reasoning~\cite{turpin2023}: the model's code says ``fetch this forbidden URL''
while its prose says ``I cannot do that.''
Which channel the operator trusts determines whether they see a compliant or
a violating system.
Raw tool-call rates are therefore not a valid safety metric without say/do
disambiguation; the same rate ($\approx 0.45$ for Sonnet, $\approx 0.32$ for
GPT-4o-mini) means opposite things across families.

\paragraph{The conformity reversal is the sharpest multi-agent finding.}
Prior LLM conformity work~\cite{weng2025, zhu2025} establishes that peer
testimony raises agreement on opinion tasks.
Our results show this does not generalize to norm-breaking:
for Sonnet, peer reports of rule-breaking make it \emph{more} compliant---it
treats social evidence of violation as a reason to scrutinize the norm, not
follow the crowd.
For GPT-4o-mini, the opposite holds: neutral peer presence alone escalates
genuine violation monotonically with group size, even when the peers explicitly
endorse following the policy (the refuser framing raises GPT-4o-mini's
violation rate at 5 peers rather than suppressing it).
In a multi-agent pipeline where agents observe each other's behavior, this
creates a risk that is invisible from single-agent evaluations: a
GPT-4o-mini--class agent that complies in isolation may escalate toward
violation as the pipeline scales up, through social facilitation rather than
explicit coordination.
The classical diffusion-of-responsibility account~\cite{darley1968} predicts
exactly this from the presence effect.

\paragraph{Honesty as a partial mitigation.}
Both a prober (Sonnet) and a genuine violator (GPT-4o-mini) accurately disclose
their logged actions to every asker, including when told the logs cannot be
verified.
This is the most practically optimistic finding: even when a model violates, it
does not appear to conceal the violation from a subsequent auditor.
The caveat is that self-reports are noisy under leading framings (GPT-4o-mini
false-confesses $\approx 50\%$ of the time on compliant transcripts), so
neutral-asker auditing is required for reliable disclosure.
Whether this honesty holds for Llama and smaller open models---which our E2/E3
runs did not cover---remains the most urgent open question, since these are the
models most likely to genuinely violate and least likely to have strong
honesty-preserving training.

\paragraph{Limitations.}
All results come from a single scenario (network-access policy, agentic
task) with a sympathetic-pull manipulation; generalization to other policy types
or task contexts is untested.
E2 and E3 for Llama were not completed.
The say/do deliver rate is estimated from small samples ($n \geq 14$) and
treated as fixed; uncertainty in that estimate propagates into all genuine-rate
calculations.
E2 used the code-execution DV throughout, so Sonnet's conformity cells measure
probing suppression rather than genuine-violation suppression, and the
informational-vs-normative belief separation called for in the pre-registered
design was not implemented.
